% ============================================================
% 分片：s05b-hecke  对应原文 PDF 第 66--76 页（印刷页 66--76）
% 负责章节：§5.3 't Hooft/Hecke 算子（含 5.3.1--5.3.7）
% 译者：trans-s05b   审校状态：初稿
% ============================================================

% ==== 原文 PDF 第 66 页（印刷页 66）====

\subsection{'t Hooft/Hecke 算子}

几何 Langlands 纲领的丰富性很大一部分来自如下事实：与局部系统对偶的 $\mathcal{D}$-模是几何 Hecke 算子的本征层。在量子场论进路中，这一点源于被称为 Wilson 算子与 't Hooft 算子的线算子之间的对偶。

我们希望描述与内窥相联系的这一图景的精细版本。为简单起见，我们聚焦于通常的例子：${}^{L}G=SO_3$，以及一个自同构群为 $\mathbb{Z}_2$ 的局部系统。

% [Tutor Note: 物理图像。全节的逻辑链是：几何 Langlands 中「Hecke 本征层」的物理化身是 't Hooft 算子的本征膜。先在电侧（${}^{L}G$ 的 Hitchin 模空间，复结构 $J$）看清 Wilson 算子 $W_p$ 的本征膜条件，再经由同调镜像对称把结论翻译为磁侧（$G$ 的 Hitchin 模空间，辛结构 $\omega_K$）的 't Hooft 算子 $T_p$ 与本征 A-膜；内窥带来的新现象恰出现在 Hitchin 纤维可约或出现轨形点的退化情形；来源：原书 p.66 上下文（A）。]

\subsubsection{Wilson 算子回顾}

首先一般地描述 Wilson 算子的作用（见 [KW] 第 8 节）。设 $(\widehat{A},\widehat{\varphi})$ 为 $SO_3$
% ==== 原文 PDF 第 67 页（印刷页 67）====
% [跨页接续: 原文第 66/67 页，句子在此跨页]
Hitchin 方程在 $\mathcal{M}_H(SO_3)\times C$ 上的万有解。于是 $\widehat{A}$ 是 $SO_3$ 丛 $\mathrm{W}\to\mathcal{M}_H(SO_3)\times C$ 上的酉联络，$\widehat{\varphi}$ 是 $\mathrm{ad}(\mathrm{W})\otimes\Omega^1_C\to\mathcal{M}_H(SO_3)\times C$ 的一个截面，使得当我们限制到某点 $m\in\mathcal{M}_H(SO_3)$ 对应的 $m\times C$ 上时，$(\widehat{A},\widehat{\varphi})$ 与由 $m$ 决定的 Hitchin 方程的解规范等价。如通常那样，$(\widehat{A},\widehat{\varphi})$ 到 $m\times C$ 的限制决定 $C$ 上的一个 $SO_3$ 局部系统，或一个 $SO_3$ Higgs 丛。

${}^{L}G$ 规范场论中的一个 Wilson 算子对应于 $C$ 上一点 $p$ 与 ${}^{L}G$ 的一个表示 ${}^{L}R$ 的选取。为简单起见，我们取 ${}^{L}R$ 为 ${}^{L}G=SO_3$ 的三维表示，并记相应的 Wilson 算子为 $W_p$。$W_p$ 在 B-膜上的作用可描述如下。设 $\mathcal{B}$ 为与凝聚层（或凝聚层复形）$\mathcal{K}\to\mathcal{M}_H$ 相联系的 B-膜。则 $W_p\cdot\mathcal{B}$ 是与层 $\mathcal{K}\otimes\mathrm{W}|_p$ 相联系的 B-膜，其中 $\mathrm{W}|_p$ 是 $\mathrm{W}$ 到 $\mathcal{M}_H\times p$ 的限制。（我们把 $\mathrm{W}$ 理解为一个秩 3 向量丛，其结构群为 $SO_3$。）因此，$W_p$ 在凝聚层上的作用为
\begin{equation}\tag{5.17}
\mathcal{K}\to\mathcal{K}\otimes\mathrm{W}|_p.
\end{equation}

% [Tutor Note: 推导补足。式 (5.17) 是本节所有电侧计算的模板：Wilson 算子并不作用在模空间上的函数，而是作用在 B-膜对应的层上，作用是「与万有丛在 $p$ 处的纤维作张量积」。这一步之所以合法，是因为 $\mathrm{W}|_p$ 对任一复结构都全纯（在 $J$ 下它是 $SO_3$ 局部系统）。由此立即可得 (5.18) 的本征条件：$\mathcal{K}\otimes\mathrm{W}|_p$ 必须形如 $\mathcal{K}\otimes V_p$，即 $\mathrm{W}|_p$ 在支撑上必须常值化；这正是下一段把光滑点上的零点膜当作电本征膜的理由；来源：原书 p.67 上下文（A）。]

此公式对构成 $\mathcal{M}_H$ 的超凯勒结构的任一复结构都有意义，因为 $\mathrm{W}$ 限制到 $\mathcal{M}_H\times p$ 后在任一复结构下都全纯。（它带有自然的联络，其曲率对每个复结构都是 $(1,1)$ 型的。）在几何 Langlands 中，人们主要关心复结构 $J$，此时 $\mathcal{M}_H(SO_3)$ 参数化 $SO_3$ 局部系统；但其他复结构的存在会简化某些计算，下文将看到这一点。

一个非常重要的特殊情形是：膜 $\mathcal{B}$ 同时是 $\mathcal{M}_H$ 的每个复结构下的 B-膜。此时我们称 $\mathcal{B}$ 为 $(B,B,B)$ 型膜——即在复结构 $I$、$J$ 或 $K$（或其任意线性组合）下的 B-膜。例子包括支撑在 $\mathcal{M}_H$ 中一点的膜——我们马上就考虑这一情形——以及由包含 $\mathcal{M}_H({}^{L}G')\subset\mathcal{M}_H({}^{L}G)$（对某个子群 ${}^{L}G'\subset{}^{L}G$）定义的膜。在后一情形，$W_r\cdot\mathcal{B}$ 仍是 $(B,B,B)$ 型膜。
% [存疑: 此处与下一句的下标均作 $W_r$，但前文只定义了 $W_p$，且所举例（点支撑膜与子群嵌入膜）与 $p$ 以外的下标无关；疑为 $W_p$ 之排版笔误，译文照原文保留 $W_r$；来源：原书 p.67 英文原文；存疑待核（A）。]
因此，$W_r$ 的作用保持 $(B,B,B)$ 型的全部拓扑对称性（即任一复结构中 B-模型的对称性）。

% [Tutor Note: 易错预警。$(B,B,B)$ 型不是「B-膜的三次方」，而是指同一对象在复结构 $I$、$J$、$K$ 下都是 B-膜；与之相对，§5.3.4 的 $(B,A,A)$ 型指在 $I$ 下是 B-膜、在 $J$ 与 $K$ 下（分别关于 $\omega_J$、$\omega_K$）是 A-膜。镜像对称交换 B-模型与 A-模型，因此电侧的 $(B,B,B)$ 对偶于磁侧的 $(B,A,A)$。判断膜的型之前必须先固定复结构或辛结构，$I$ 与 $\Omega_I$、$J$ 与 $\omega_J$、$K$ 与 $\omega_K$ 三组约定不可混用；来源：原书 p.67 上下文（A）。]

特别令人感兴趣的是 Wilson 算子的本征膜，也称为电本征膜。膜与向量空间 $V$ 的张量积定义如下：若 $\mathcal{B}$ 由层 $\mathcal{K}$ 定义，则 $\mathcal{B}\otimes V$ 由层 $\mathcal{K}\otimes V$ 定义。若
\begin{equation}\tag{5.18}
W_p\cdot\mathcal{B}=\mathcal{B}\otimes V_p,
\end{equation}
对某个向量空间 $V_p$ 成立，则称膜 $\mathcal{B}$ 为 $W_p$ 的本征膜。我们称 $V_p$ 为乘子（multiplier）。在复结构 $J$ 下，$\mathcal{M}_H$ 的 B-模型与一个四维拓扑场论相关，利用一般性论证（见 [KW] 第 6.4 节）可以证明：若式 (5.18) 对 $p$ 的一个取值成立，则它对所有 $p$ 成立，且 $V_p$ 作为某个局部系统的纤维随 $p$ 变化。

与 (5.17) 比较可知，要使膜 $\mathcal{B}$ 成为电本征膜，丛 $\mathrm{W}|_p$ 在限制到 $\mathcal{B}$ 的支撑时必须平凡——等价于一个常向量空间。若 $\mathcal{B}$ 的支撑是 $\mathcal{M}_H(SO_3)$ 中的光滑点 $m$，则情形正是如此。更具体地说，我们取 $\mathcal{B}$ 为与支撑在 $m$ 的点层相联系的膜（物理学文献中称为零点膜（zero-brane））。这样的膜是电本征膜，其乘子为向量空间 $\mathrm{W}|_{m\times p}$，即
% ==== 原文 PDF 第 68 页（印刷页 68）====
% [跨页接续: 原文第 67/68 页，句子在此跨页]
$\mathrm{W}\to\mathcal{M}_H\times C$ 到点 $m\times p\in\mathcal{M}_H\times C$ 的限制：
\begin{equation}\tag{5.19}
W_p\cdot\mathcal{B}=\mathcal{B}\otimes\mathrm{W}|_{m\times p}.
\end{equation}

此陈述在任一复结构下都成立，因此我们可以把 $\mathcal{B}$ 看作 $(B,B,B)$ 型本征膜。换言之，它是 $\mathcal{M}_H$ 在复结构 $I$、$J$ 或 $K$（或其任意组合）下的 B-膜，并且在任一复结构下都是本征膜。在几何 Langlands 纲领中，人们主要关心复结构 $J$；但零点膜同时在三个复结构下都是本征膜这一事实会简化计算，这一点将逐渐清楚。

\subsubsection{轨形奇点处的 B-膜}

我们希望把上述分析重复到支撑在 $\mathbb{Z}_2$ 轨形奇点 $r\in\mathcal{M}_H$ 的膜的情形。这样的奇点对应于一个结构群约化到 $O_2$ 的 $SO_3$ 局部系统。回忆 $O_2$ 作为子群嵌入 $SO_3$：
\begin{equation}\tag{5.20}
\begin{pmatrix}
* & * & 0\\
* & * & 0\\
0 & 0 & \pm 1
\end{pmatrix}
\end{equation}
并且任何结构群约化到 $O_2$ 的 $SO_3$ 局部系统都以 $\mathbb{Z}_2$ 为对称群，由 $O_2$ 的中心元素生成：
\begin{equation}\tag{5.21}
\begin{pmatrix}
-1 & 0 & 0\\
0 & -1 & 0\\
0 & 0 & 1
\end{pmatrix}.
\end{equation}

如通常那样，我们考虑这种类型的一般局部系统，其自同构群恰为这个 $\mathbb{Z}_2$。在当前语境下，$\mathrm{W}|_r$（即 $\mathrm{W}$ 到 $r\times C$ 的限制）是一个结构群约化到 $O_2$ 的局部系统，因此有分解
\begin{equation}\tag{5.22}
\mathrm{W}|_r=U\oplus\det U,
\end{equation}
其中 $U$ 是一个秩 2 局部系统，结构群为 $O_2$，$\det U$ 是它的行列式。$\mathbb{Z}_2$ 的中心生成元在 $U$ 上作用为 $-1$，在 $\det U$ 上作用为 $+1$。

正如 §2.2 所解释的，支撑在轨形奇点 $r$ 的膜范畴由两个不可约对象 $\mathcal{B}_+$ 与 $\mathcal{B}_-$ 生成。它们各自与支撑在 $r$ 的点层相联系；二者的区别在于 $\mathbb{Z}_2$ 的非平凡元素在该层上作用为乘以 $+1$ 还是乘以 $-1$。

当我们用 Wilson 算子 $W_p$ 作用于 $\mathcal{B}_+$ 或 $\mathcal{B}_-$ 时会发生什么？既然 $\mathcal{B}_+$ 与 $\mathcal{B}_-$ 都以点层支撑于 $r$，$W_p$ 对它们的作用便是与三维向量空间 $\mathrm{W}|_{r\times p}$（$\mathrm{W}$ 在 $r\times p$ 处的纤维）作张量积。不过，为追踪 $\mathbb{Z}_2$ 作用，我们应更精确些。由式 (5.22)，有分解 $\mathrm{W}|_{r\times p}=U|_p\oplus\det U|_p$，其中 $\mathbb{Z}_2$ 的非平凡元素在第一直和项上作用为 $-1$、在第二直和项上作用为 $+1$。于是
\begin{equation}\tag{5.23}
\begin{aligned}
W_p\cdot\mathcal{B}_+ &= (\mathcal{B}_-\otimes U|_p)\oplus(\mathcal{B}_+\otimes\det U|_p)\\
W_p\cdot\mathcal{B}_- &= (\mathcal{B}_+\otimes U|_p)\oplus(\mathcal{B}_-\otimes\det U|_p).
\end{aligned}
\end{equation}

% [Tutor Note: 易错预警。$\mathcal{B}_\pm$ 按 $\mathbb{Z}_2$ 的非平凡元素作用为 $+1$ 或 $-1$ 而区分。$\det U|_p$ 上的作用是平凡的，故张量积不改变 $\mathbb{Z}_2$-特征标：$\mathcal{B}_\pm\otimes\det U|_p$ 与 $\mathcal{B}_\pm$ 同构；而 $U|_p$ 上的作用为 $-1$，故 $\mathcal{B}_\pm\otimes U|_p$ 同构于两个 $\mathcal{B}_\mp$ 之和，系数 $2=\dim U$ 来自 $U$ 的秩。这些同构都非典范（依赖 $U|_p$ 的基的选取），(5.24) 只是同构意义下的陈述；不要把它误读为两个不同的不可约膜；来源：原书 p.68 上下文（A）。]

% ==== 原文 PDF 第 69 页（印刷页 69）====

一个不那么精确但更具启发性的改写方式如下。$\det U|_p$ 是 $\mathbb{Z}_2$ 平凡作用的一维向量空间，故 $\mathcal{B}_\pm\otimes\det U|_p$ 非典范地同构于 $\mathcal{B}_\pm$。而 $U|_p$ 是 $\mathbb{Z}_2$ 的非平凡元素作用为乘以 $-1$ 的二维向量空间，故 $\mathcal{B}_\pm\otimes U|_p$ 非典范地同构于两个 $\mathcal{B}_\mp$ 之和。于是，在同构意义下有
\begin{equation}\tag{5.24}
\begin{aligned}
W_p\cdot\mathcal{B}_+ &= \mathcal{B}_+ + 2\mathcal{B}_-\\
W_p\cdot\mathcal{B}_- &= \mathcal{B}_- + 2\mathcal{B}_+.
\end{aligned}
\end{equation}

我们希望理解这些陈述的磁对偶。在 §5.3.3 与 §5.3.4 中，我们回顾一般 Hitchin 纤维上的几何 Langlands 对偶；然后在 §5.3.5 中考虑与内窥相关的特殊 Hitchin 纤维的情形。

\subsubsection{\texorpdfstring{$\varphi$}{φ}-不变的 Hecke 修改}

Wilson 算子 $W_p$ 的磁对偶是 't Hooft 算子 $T_p$。关于 't Hooft 算子的定义及其与几何 Langlands 纲领中通常的 Hecke 算子的关系，见 [KW] 第 9、10 节。若对每个 't Hooft 算子 $T_p$ 都有
\begin{equation}\tag{5.25}
T_p\cdot\mathcal{A}=\mathcal{A}\otimes V_p
\end{equation}
对某个向量空间 $V_p$ 成立，则称 A-膜 $\mathcal{A}$ 为 't Hooft 算子的本征膜；这样的膜称为磁本征膜。${}^{L}G$ 规范场论的 Wilson 算子按 ${}^{L}G$ 的表示分类，$G$ 规范场论的 't Hooft 算子同样按 ${}^{L}G$ 的表示分类。电-磁对偶预期把 Wilson 算子映为 't Hooft 算子、把电本征膜映为磁本征膜。

下面回顾 't Hooft 算子 $T_p$ 对 Higgs 丛 $(E,\varphi)$ 的作用。当 $\varphi=0$ 时，可能的 Hecke 修改就是几何 Langlands 纲领中通常考虑的那些；它们由仿射 Grassmann 簇的一个子簇（称为 Schubert 簇（Schubert variety））$\mathcal{V}$ 参数化，$\mathcal{V}$ 依赖于对偶群 ${}^{L}G$ 的一个表示 ${}^{L}R$ 的选取。例如，继续用我们的例子：若 $G=SL_2$ 而 ${}^{L}R$ 是 ${}^{L}G=SO_3$ 的三维表示，则 $\mathcal{V}$ 中的一般点对应 $SL_2$ 丛 $E$ 的如下 Hecke 修改：对 $E$ 作为线丛之和的某个局部分解 $E=\mathcal{N}_1\oplus\mathcal{N}_2$，$E$ 被映为 $\mathcal{N}_1(p)\oplus\mathcal{N}_2(-p)$。让 $\mathcal{N}_1$ 与 $\mathcal{N}_2$ 变化，就得到 $E$ 的两参数族 Hecke 修改。$E$ 的这种类型的修改族可以退化到平凡修改，而 $\mathcal{V}$ 含有对应平凡 Hecke 变换的点。

% [Tutor Note: 易错预警。注意 Hecke 修改 $E\mapsto\mathcal{N}_1(p)\oplus\mathcal{N}_2(-p)$ 中 $+p$ 与 $-p$ 成对出现，故 $\det E$ 不变：$\det(\mathcal{N}_1(p)\oplus\mathcal{N}_2(-p))=\det E$。这是 $SL_2$（而非 $GL_2$）的 Hecke 修改必须满足的条件。可与 §5.3.6 相对照：那里 $SO_3$ 侧的算子通过 $E\mapsto\mathcal{N}_1(p)\oplus\mathcal{N}_2$ 作用，恰好改变 $\det E$ 的次数的模 2 约化，这正是两侧算子对偶方式不同的体现；若在上面漏掉 $-p$，所构造的就不是 $SL_2$ 丛之间的修改；来源：原书 p.69 上下文（A）。]

我们刚刚对这个例子所描述的，是 't Hooft 算子在 $\varphi=0$ 这一最退化情形下可能的作用。反之，若 $\varphi\neq0$，则必须把限制在某些意义下 $\varphi$-不变的 Hecke 修改上。对 $G=SL_2$，并假设 $\varphi$ 在点 $p$ 处正则半单，$\varphi$-不变性意味着分解 $E=\mathcal{N}_1\oplus\mathcal{N}_2$ 必须与 $\varphi$ 的作用相容，即 $\varphi:E\to E\otimes K$ 把 $\mathcal{N}_1$ 映到 $\mathcal{N}_1\otimes K$、把 $\mathcal{N}_2$ 映到 $\mathcal{N}_2\otimes K$。$\mathcal{N}_1$ 与 $\mathcal{N}_2$ 的可能选取恰有两种：局部地，由于 $\varphi(p)$ 正则半单，我们可以对角化 $\varphi$
\begin{equation}\tag{5.26}
\varphi=\begin{pmatrix}
a & 0\\
0 & -a
\end{pmatrix},
\end{equation}
而 $\mathcal{N}_1$ 与 $\mathcal{N}_2$ 必须（不计排列）等于两个「特征子空间」。

% [Tutor Note: 推导补足。$\varphi(p)$ 正则半单等价于两个特征值互异且非零，即 (5.26) 中 $a\neq0$；此时特征子空间都是一维的，故 $\mathcal{N}_1$、$\mathcal{N}_2$ 作为线丛被唯一确定（不计次序）。原文的「up to permutation」对应如下事实：$p$ 上方两点 $p'$、$p''$ 没有自然的标号，交换标号即交换 $\mathcal{N}_1$ 与 $\mathcal{N}_2$，这与后文 $\Phi$、$\Phi^{-1}$ 的互换以及关于 $p$ 的单值作用一致；来源：原书 p.69 上下文（A）。]

% ==== 原文 PDF 第 70 页（印刷页 70）====

除了这两个非平凡的 $\varphi$-不变 Hecke 修改之外，还必须记得计入平凡的 Hecke 修改（因为它对应 $\mathcal{V}$ 中的一个点），它也是 $\varphi$-不变的。在这个例子中总共有三个 $\varphi$-不变 Hecke 修改，即一个平凡的与两个非平凡的；这一陈述（我们将会看到）对偶于我们出发时 ${}^{L}G=SO_3$ 的表示是三维的这一事实。

现在让我们从谱曲线 $\pi:D\to C$ 的角度来看 $\varphi$-不变 Hecke 修改是什么样子。先考虑一般谱曲线的情形，它由方程 $\det(z-\varphi)=0$ 给出。$\varphi$-不变 Hecke 修改保持 $\varphi$ 的特征多项式不变，从而把 Hitchin 纤维化的每个纤维 $\boldsymbol{F}$ 映到自身。它如何作用在 $\boldsymbol{F}$ 上？

使 $\varphi$ 正则半单的点 $p\in C$ 上方有两个不同的点 $p',p''\in D$。丛 $E$ 是 $\pi_*(\mathcal{L})$（对某个线丛 $\mathcal{L}\to D$），且 $\varphi=\pi_*(z)$。后一条件意味着 $\varphi(p)$ 的特征子空间对应 $p$ 上方 $z$ 的两个不同取值，换言之对应两个点 $p'$ 与 $p''$。既然如此，$(E,\varphi)$ 在点 $p$ 处的非平凡 $\varphi$-不变 Hecke 修改来自 $\mathcal{L}$ 的如下特定形式的变换
\begin{equation}\tag{5.27}
\mathcal{L}\to\mathcal{L}\otimes\mathcal{O}(p'-p''),
\end{equation}
对 $p$ 上方两点 $p',p''$ 的两种可能标号之一而言。（Higgs 丛 $(E,\varphi)$ 的 Hecke 修改这一概念在数学上是自然的，并被 [DP] 取作出发点。）把它下推到 $C$ 时，它以所希望的方式修改 $E$。

% [Tutor Note: 推导补足与易错预警。记号 $\varphi=\pi_*(z)$ 是简写：$z$ 是谱曲线 $D$（由 $\det(z-\varphi)=0$ 切出，位于 $K_C$ 的全空间内）上的坐标函数，$\varphi$ 是「乘以 $z$」这一映射 $\mathcal{L}\to\mathcal{L}\otimes\pi^*K_C$ 的推前。原文用星号位置区分推前与拉回：本页 $\pi_*(\mathcal{L})$ 的星号在下方（推前），而 §5.3.6 中 $\pi^*(\mathcal{O}(-p))$ 的星号在上方（拉回），排印上极易混淆，引用时务必核对；来源：原书 p.70 上下文（A）。]

现在我们可以看出为什么支撑在 Hitchin 纤维化的纤维 $\boldsymbol{F}$ 上、并带有平坦线丛 $\mathcal{R}$ 的 A-膜 $\mathcal{A}_{\boldsymbol{F}}$ 是磁本征膜，即 't Hooft 算子 $T_p$ 的本征膜。首先，$T_p$ 把 $\boldsymbol{F}$ 映到自身，因为它保持 $\varphi$ 的特征多项式。既然 $T_p$ 保持 $\mathcal{A}_{\boldsymbol{F}}$ 的支撑，$\mathcal{A}_{\boldsymbol{F}}$ 成为 $T_p$ 的本征膜就是可以设想的。

现在，假设我们选取 $p$ 使 $\varphi(p)$ 正则半单（我们只处理这一情形），则 $T_p\cdot\mathcal{A}_{\boldsymbol{F}}$ 的赋值来自刚才描述的三个 $\varphi$-不变 Hecke 修改的贡献之和。其中一个是平凡 Hecke 修改，它使 $\mathcal{A}_{\boldsymbol{F}}$ 不变。另外两个来自变换 $\mathcal{L}\to\mathcal{L}\otimes\mathcal{O}(p'-p'')$（对两点的某种标号）。这样的变换可以解释为 Hitchin 纤维的同构 $\Phi:\boldsymbol{F}\to\boldsymbol{F}$。若两点 $p'$ 与 $p''$ 的标号反过来，则 $\Phi$ 被替换为 $\Phi^{-1}$。

$\boldsymbol{F}$ 是复环面，$\Phi$ 是 $\boldsymbol{F}$ 沿常向量的「平移」。一般地，若 $\mathcal{R}\to\boldsymbol{F}$ 是复环面上的平坦线丛、$\Phi:\boldsymbol{F}\to\boldsymbol{F}$ 是平移，则 $\Phi^*(\mathcal{R})=\mathcal{R}\otimes\mathcal{V}$，其中 $\mathcal{V}$ 是某个一维向量空间。由此可知 $\mathcal{A}$ 是 $T_p$ 的本征膜。事实上我们有
\begin{equation}\tag{5.28}
T_p\cdot\mathcal{A}_{\boldsymbol{F}}=\mathcal{A}_{\boldsymbol{F}}\otimes(\mathbb{C}\oplus\mathcal{V}\oplus\mathcal{V}^{-1}),
\end{equation}
其中右端三项分别来自平凡 Hecke 修改，以及涉及 $\Phi$ 与 $\Phi^{-1}$ 的非平凡修改。

% [Tutor Note: 推导补足与物理图像。$\Phi$ 是复环面 $\boldsymbol{F}$ 的平移，而平移下平坦线丛只差一个一维向量空间：$\Phi^*(\mathcal{R})=\mathcal{R}\otimes\mathcal{V}$，故三个贡献给出乘子 $\mathbb{C}\oplus\mathcal{V}\oplus\mathcal{V}^{-1}$：$\mathbb{C}$ 来自平凡 Hecke 修改（对偶于电侧 $\varphi(p)$ 的核），$\mathcal{V}$ 与 $\mathcal{V}^{-1}$ 分别来自 $\Phi$ 与 $\Phi^{-1}$。这与电侧 $E_p$ 按 $\varphi(p)$ 分解为「核 $\oplus$ 两个特征子空间」逐项对偶，是本节最值得逐项核对的对偶关系；来源：原书 p.70 上下文（A）。]

让我们把这与电侧的情形相比较。那里我们考虑的膜 $\mathcal{B}$，其支撑在 $\mathcal{M}_H(SO_3)$ 中对应一个 Higgs 丛 $(E,\varphi)$。它满足
% ==== 原文 PDF 第 71 页（印刷页 71）====
% [跨页接续: 原文第 70/71 页，句子与公式在此跨页]
$W_p\cdot\mathcal{B}=\mathcal{B}\otimes E_p$，其中 $E_p$ 是 $E$ 在 $p$ 处的纤维。由于我们假设 $\varphi(p)$ 正则半单（这是在对偶下不变的性质），可以按 $\varphi(p)$ 的特征子空间分解 $E_p$。$\varphi(p)$ 有两个非零特征值的特征子空间；由于 $E_p$ 上的二次型，它们互为对偶，故我们记它们为 $\mathcal{X}$ 与 $\mathcal{X}^{-1}$。它们与 $\mathcal{V}$ 和 $\mathcal{V}^{-1}$ 有自然的同构；这可用与证明 $GL_1$ 的几何 Langlands 对偶时类似的办法建立。关于这一计算的一种版本见 [DP]。$\varphi(p)$ 的核对应直和项 $\mathbb{C}$。

\subsubsection{解释}

也许最需要澄清的是对刚刚描述的结果的解释。

那个计算使用了 Hitchin 纤维化以及适用于复结构 $I$ 下 B-模型的其他工具。然而对于几何 Langlands 纲领，我们真正关心的是辛结构 $\omega_K$ 下 A-模型的 A-膜。

我们前面已注意到，支撑在一点的电本征膜是 $(B,B,B)$ 型膜——即它在复结构 $I$、$J$、$K$ 的每一个下都是 B-膜。对偶的陈述是：对偶的磁本征膜是 $(B,A,A)$ 型膜。$(B,A,A)$ 型膜是这样一个膜：它同时是复结构 $I$ 下的 B-膜，以及辛结构 $\omega_J$ 与 $\omega_K$ 的 A-模型的 A-膜。

为显式看出支撑在 Hitchin 纤维化的纤维 $\boldsymbol{F}$ 上、并带有酉平坦线丛 $\mathcal{L}$ 的膜 $\mathcal{A}_{\boldsymbol{F}}$ 是 $(B,A,A)$ 型，我们论证如下。$\boldsymbol{F}$ 是复结构 $I$ 下的复拉格朗日子流形。$\mathcal{A}_{\boldsymbol{F}}$ 是复结构 $I$ 下的 B-膜，因为 $\boldsymbol{F}$ 在该复结构下全纯，而像 $\mathcal{L}$ 这样的平坦线丛也全纯。$\boldsymbol{F}$ 关于 $\Omega_I=\omega_J+i\omega_K$ 是拉格朗日的，从而关于 $\omega_J$ 与 $\omega_K$ 也是拉格朗日的。最标准的 A-膜类型是带酉平坦线丛的拉格朗日子流形；无论把辛结构取为 $\omega_J$ 还是 $\omega_K$，$\mathcal{A}_{\boldsymbol{F}}$ 都符合这一类型。

Wilson 算子 $W_p$ 保持 $(B,B,B)$ 型超对称这一事实的对偶是：'t Hooft 算子 $T_p$ 保持 $(B,A,A)$ 型超对称。这一陈述意味着：若 $\widetilde{\mathcal{A}}$ 是 $(B,A,A)$ 型膜，则 $T_p\cdot\widetilde{\mathcal{A}}$ 也是。（更一般地，若 $\widetilde{\mathcal{A}}$ 保持这些超对称的任一线性组合，则 $T_p\cdot\widetilde{\mathcal{A}}$ 亦然。）要真正确定 $T_p\cdot\widetilde{\mathcal{A}}$ 是哪一个 $(B,A,A)$ 型膜，我们可以用最方便的那一个超对称来计算。代数几何是有力的工具，因此把 $T_p\cdot\widetilde{\mathcal{A}}$ 视为复结构 $I$ 下的 B-膜来计算很可能是方便的。这足以把 $T_p\cdot\widetilde{\mathcal{A}}$ 确定为 $(B,A,A)$ 型膜，特别是确定为 $K$ 型的 A-膜，而这正是几何 Langlands 所需。

最后一句之所以成立，是因为给定一个 $I$ 型 B-膜，至多只有一种方式赋予它 $(B,A,A)$ 型膜的结构。让我们在当前语境下具体说明。§5.3.3 中我们用谱曲线与代数几何，以全纯子簇 $\boldsymbol{F}\subset\mathcal{M}_H(SL_2)$ 与全纯向量丛 $\mathcal{K}\to\boldsymbol{F}$ 的数据构造了膜 $T_p\cdot\mathcal{A}_{\boldsymbol{F}}$。这些数据决定复结构 $I$ 下的一个 B-膜。赋予该膜 $(B,A,A)$ 型结构，意味着赋予 $\mathcal{K}$ 一个埃尔米特度量，使诱导的酉联络平坦；带上这样平坦丛的 $\boldsymbol{F}$ 就是 $(B,A,A)$ 型膜。
% ==== 原文 PDF 第 72 页（印刷页 72）====
% [跨页接续: 原文第 71/72 页，句子在此跨页]
使全纯丛 $\mathcal{K}$ 可被如此赋予的平坦酉联络，若存在，则在同构意义下唯一。在当前语境下，由于 $T_p$ 保持 $(B,A,A)$ 型超对称，我们预期这样的平坦联络存在，而 §5.3.3 的计算表明它确实存在。

最后一点评论涉及公式中表达 $\mathcal{A}_{\boldsymbol{F}}$ 是磁本征膜这一事实的「乘子」$V_p$：
\begin{equation}\tag{5.29}
T_p\cdot\mathcal{A}_{\boldsymbol{F}}=\mathcal{A}_{\boldsymbol{F}}\otimes V_p.
\end{equation}
我们希望讨论 $p$ 变化时会发生什么。在几何 Langlands 中，我们预期 $V_p$ 是某个局部系统的纤维，即一个平坦向量丛。这并不是我们从基于谱曲线与复结构 $I$ 中复几何的描述所能最直接得到的东西。从那个观点看，我们得到的 $V_p$ 是一个秩 3 全纯向量丛 $V\to C$ 的纤维（该丛起初只在 $\varphi$ 正则半单处有定义），而它没有明显的平坦结构。然而，乘以谱曲线的坐标 $z$ 给出一个 $K$-值自同态，我们记为 $\theta$，因此乘子实际上是一个 Higgs 丛 $(V,\theta)$。通过解 Hitchin 方程，我们把一个秩 3 局部系统与这个 Higgs 丛联系起来，而我们预期它与我们出发时的 $SO_3$ 局部系统同构。

\subsubsection{可约纤维}

现在我们准备考虑与内窥有关的情形。

我们考虑 Hitchin 纤维化的一个特殊纤维 $\boldsymbol{F}$，它是两个不可约分支 $\boldsymbol{F}_1$ 与 $\boldsymbol{F}_2$ 的并，且这两分支彼此在一个除子上相交。既然如此，我们可以构造支撑在 $\boldsymbol{F}_1$ 或 $\boldsymbol{F}_2$ 上的秩 1 A-膜 $\mathcal{A}_1$ 与 $\mathcal{A}_2$。若 $\boldsymbol{F}_1$ 与 $\boldsymbol{F}_2$ 单连通，这些膜是唯一的（如亏格 1、一个分歧点的曲线的情形）；否则它们依赖于我们未显式指出的参数。

在推导描述 $T_p$ 作用的式 (5.28) 时，一个关键要素是由 $\mathcal{L}\to\mathcal{L}\otimes\mathcal{O}(p'-p'')$ 给出的映射 $\Phi:\boldsymbol{F}\to\boldsymbol{F}$。$\boldsymbol{F}$ 可约时本质的新事实仅仅是：$\Phi$ 交换 $\boldsymbol{F}$ 的两个分支。这正是我们在 §5.2.1 中刻画两个分支的方式。同样地，$\Phi^{-1}$ 交换两个分支。因此 $\Phi$ 或 $\Phi^{-1}$ 交换 $\mathcal{A}_1$ 与 $\mathcal{A}_2$。既然 $T_p$ 通过 $1+\Phi+\Phi^{-1}$ 作用，便可得在同构意义下
\begin{equation}\tag{5.30}
\begin{aligned}
T_p\cdot\mathcal{A}_1 &= \mathcal{A}_1+2\mathcal{A}_2\\
T_p\cdot\mathcal{A}_2 &= \mathcal{A}_2+2\mathcal{A}_1.
\end{aligned}
\end{equation}
这与电侧公式 (5.24) 完全平行。

若 $\mathcal{A}_1$ 与 $\mathcal{A}_2$ 有模，则应稍更精确地描述。$\mathcal{A}_1$ 依赖于适当线丛 $\mathcal{L}\to\boldsymbol{F}_1$ 的选取，且我们应取 $\mathcal{A}_2$ 为与线丛 $\Phi^*(\mathcal{L})\to\boldsymbol{F}_2$ 相联系的膜。注意 $\Phi^*(\mathcal{L})$ 与 $(\Phi^{-1})^*(\mathcal{L})$ 同构，但并非典范同构。

预期可通过 §5.3.4 的步骤得到与 (5.23) 类似的更精确结果。人们用代数几何的标准方法，构造 $T_p\cdot\mathcal{A}_1$ 与 $T_p\cdot\mathcal{A}_2$ 作为复结构 $I$ 下的 B-膜。这将给出比
% ==== 原文 PDF 第 73 页（印刷页 73）====
% [跨页接续: 原文第 72/73 页，句子在此跨页]
(5.30) 更精确的结果：
\begin{equation}\tag{5.31}
\begin{aligned}
T_p\cdot\mathcal{A}_1 &= (\mathcal{A}_1\otimes\mathcal{J}_1)\oplus(\mathcal{A}_2\otimes\mathcal{J}_2)\\
T_p\cdot\mathcal{A}_2 &= (\mathcal{A}_2\otimes\mathcal{J}_1')\oplus(\mathcal{A}_1\otimes\mathcal{J}_2'),
\end{aligned}
\end{equation}
其中 $\mathcal{J}_1,\mathcal{J}_2$ 等是向量空间，其维数由下标指示。它们都带有自然的 $K$-值自同态 $\theta_1,\theta_2$ 等，来自 Higgs 场（即乘以谱曲线的坐标 $z$）；且 $(\mathcal{J}_1,\theta_1)$ 等是 $C$ 上的 Higgs 丛。把这些 Higgs 丛经 Hitchin 方程与局部系统联系起来，预期可得到 (5.23) 的类比：
\begin{equation}\tag{5.32}
\begin{aligned}
T_p\cdot\mathcal{A}_1 &= (\mathcal{A}_2\otimes U|_p)\oplus(\mathcal{A}_1\otimes\det U|_p)\\
T_p\cdot\mathcal{A}_2 &= (\mathcal{A}_1\otimes U|_p)\oplus(\mathcal{A}_2\otimes\det U|_p).
\end{aligned}
\end{equation}

% [Tutor Note: 推导补足。(5.31) 是代数几何计算的中间结果：$\mathcal{J}_i$ 只是带 $K$-值自同态 $\theta_i$ 的向量空间（即 $C$ 上 Higgs 丛 $(\mathcal{J}_i,\theta_i)$ 的纤维），其下标指示维数；用 Hitchin 方程把 $(\mathcal{J}_i,\theta_i)$ 换成局部系统后，才得到与 (5.23) 形式相同的 (5.32)。阅读时不要把 (5.31) 误当作最终的对偶陈述；来源：原书 p.73 上下文（A）。]

\subsubsection{互反情形}

我们可以把类似的技术用于互反情形 ${}^{L}G=SL_2$、$G=SO_3$。

对规范群 $SL_2$，要考虑的基本 Wilson 算子是二维表示对应的算子 $\widetilde{W}_p$。粗略地说，它通过式 (5.17) 的明显类比作用。设 $(\mathrm{E},\widetilde{\varphi})$ 表示 $\mathcal{M}_H(SL_2)\times C$ 上的万有 Higgs 丛，则 $\widetilde{W}_p$ 对定义 B-膜 $\mathcal{B}$ 的层 $\mathcal{K}$ 的作用为
\begin{equation}\tag{5.33}
\mathcal{K}\to\mathcal{K}\otimes\mathrm{E}|_p,
\end{equation}
其中 $\mathrm{E}|_p$ 是万有秩 2 丛 $\mathrm{E}\to\mathcal{M}_H\times C$ 到 $\mathcal{M}_H\times p$ 的限制。$\widetilde{W}_p$ 满足
\begin{equation}\tag{5.34}
\widetilde{W}_p^2=1+W_p,
\end{equation}
这表达了如下事实：二维表示与自身的张量积是平凡表示与三维表示的直和；它们对应式 (5.34) 右端的项 $1$ 与 $W_p$。

一个重要的微妙之处反映了 $SL_2$ 的中心在二维表示中非平凡地作用这一事实。万有丛 $\mathrm{E}$ 作为通常意义下的向量丛并不存在。它必须被理解为扭向量丛，被 $\mathcal{M}_H$ 上某个 $\mathbb{C}^\times$ gerbe 所扭。所涉及的 gerbe 由 $\mathbb{Z}_2$ gerbe（$\mathbb{Z}_2$ 在此是 $SL_2$ 的中心）经包含 $\mathbb{Z}_2\subset\mathbb{C}^\times$ 诱导而来。由于 $\mathrm{E}|_p$ 是扭向量丛，与它作张量积把 $\mathcal{M}_H$ 上的普通层映为扭层，反之亦然。用 [KW] 的语言说，这意味着 $\widetilde{W}_p$ 在膜上的作用平移离散电通量 $\boldsymbol{e}_0$，后者是 $SL_2$ 中心的一个特征标。对偶的陈述是：对偶 't Hooft 算子 $\widetilde{T}_p$ 平移离散磁通量 $\boldsymbol{m}_0$，后者是一个 $SO_3$ 丛的第二 Stiefel--Whitney 类 $w_2$。
% [原文疑误: 原书 p.73 此处将 Stiefel 误拼为 "Stieffel"（多一个 f）；译文按正确拼写译为 Stiefel--Whitney，判断依据是图像 p.73 第 17 行与标准术语；来源：原书 p.73 英文原文与译本对照（A）。]

% [Tutor Note: 易错预警与物理图像。$\mathrm{E}$ 只是扭向量丛（被 $\mathbb{C}^\times$ gerbe 扭），故与 $\mathrm{E}|_p$ 作张量积会把普通层与扭层互换：$\widetilde{W}_p$ 把 $\mathcal{B}$ 送到 $\mathcal{B}'$、把 $\mathcal{B}'$ 送到 $\mathcal{B}$，单看任何一个都不是从同一范畴到自身的算子，只有直和 $\widehat{\mathcal{B}}$ 才是通常意义的本征膜。取平方后张量积不再改变扭化，这正对应 (5.34) 的分解 $2\otimes2=1\oplus3$：右端的 $1$ 与 $W_p$ 都是 $SO_3$ 侧合法的算子。相应的守恒量是离散电通量 $\boldsymbol{e}_0$（$SL_2$ 中心的特征标）与离散磁通量 $\boldsymbol{m}_0=w_2$；$\widetilde{W}_p$ 平移前者，$\widetilde{T}_p$ 平移后者；来源：原书 p.73 上下文（A）。]

粗略地说，支撑在光滑点 $r\in\mathcal{M}_H(SL_2)$ 的点层给出电本征膜 $\mathcal{B}$，正如我们先前对 ${}^{L}G=SO_3$ 的讨论。但这略有简化。支撑在点 $r$ 的点层既可以理解为普通层，也可以理解为扭层，因为扭化涉及的 gerbe 在限制到光滑点时平凡。与 $\mathrm{E}$ 作张量积会把
% ==== 原文 PDF 第 74 页（印刷页 74）====
% [跨页接续: 原文第 73/74 页，句子在此跨页]
这两种情形互换。因此，若用 $\mathcal{B}$ 与 $\mathcal{B}'$ 分别表示与点层相联系的膜的普通版本与扭版本，则 Wilson 算子的作用为
\begin{equation}\tag{5.35}
\begin{aligned}
\widetilde{W}_p\cdot\mathcal{B} &= \mathcal{B}'\otimes\mathrm{E}|_{r\times p}\\
\widetilde{W}_p\cdot\mathcal{B}' &= \mathcal{B}\otimes\mathrm{E}|_{r\times p}.
\end{aligned}
\end{equation}
因此，和 $\widehat{\mathcal{B}}=\mathcal{B}\oplus\mathcal{B}'$ 在通常意义下是电本征膜：
\begin{equation}\tag{5.36}
\widetilde{W}_p\cdot\widehat{\mathcal{B}}=\widehat{B}\otimes\mathrm{E}|_{r\times p}.
\end{equation}
% [原文疑误: (5.36) 右端作斜体 $B$ 上的 $\widehat{B}$，而左端及紧邻上文均作花体 $\widehat{\mathcal{B}}$；二者应为同一对象，此处字体不一致，疑为排版笔误，建议统一为花体；来源：原书 p.74 英文原文与译本对照（A）。]

对偶 't Hooft 算子 $\widetilde{T}_p$ 在 $SO_3$ 丛 $W$ 或 $SO_3$ Higgs 丛 $(W,\varphi)$ 上的作用与 §5.3.3 中所描述的相当类似。用使得 $W=\mathrm{ad}(E)$ 的 $SL_2$ Higgs 丛 $(E,\varphi)$ 来描述这一作用是方便的，其中 $\det E$ 的次数模 2 同余于 $w_2(W)$。$E$ 并非完全唯一确定，但以下陈述用 $W=\mathrm{ad}(E)$ 表达时不依赖于 $E$ 的选取。相对于 $E$ 作为线丛之和的某个局部分解 $E=\mathcal{N}_1\oplus\mathcal{N}_2$，$\widetilde{T}_p$ 通过 $E\to\mathcal{N}_1(p)\oplus\mathcal{N}_2$ 作用。这一运算反转 $\det E$ 的次数模 2 的约化，因此如所预期地反转 $\boldsymbol{m}_0$。

若 $\varphi=0$，选取 $\mathcal{N}_1$ 的自由导致由 $\mathbb{CP}^1$ 参数化的一族可能的 Hecke 修改。这个参数空间是紧的（反映 $SL_2$ 的二维表示是 minuscule（极小）表示这一事实），故无需再添加任何东西使之紧化。若 $\varphi\neq0$，$\widetilde{T}_p$ 通过 $\varphi$-不变的 Hecke 修改作用。若 $\varphi(p)$ 正则半单，这意味着在上一段中我们必须取 $\mathcal{N}_1$ 为 $\varphi$ 的两个特征子空间之一。两种选取的存在对偶于我们出发时 $SL_2$ 的表示是二维的这一事实。

用谱曲线来说，$\varphi$-不变 Hecke 修改的作用可以非常类似式 (5.27) 地描述。使 $E=\pi_*(\mathcal{L})$ 的线丛 $\mathcal{L}\to D$ 按
\begin{equation}\tag{5.37}
\mathcal{L}\to\mathcal{L}(p^*)
\end{equation}
变换，其中 $p^*$ 可以是 $p$ 上方两点 $p'$ 与 $p''$ 中的任一个。

现在 $\mathcal{M}_H(SO_3)$ 有两个分支，按 $w_2(W)$ 分类。记 $\boldsymbol{F}$ 与 $\boldsymbol{F}_\theta$ 为这两个分支的 Hitchin 纤维化的纤维。当 $p^*$ 等于 $p'$ 或 $p''$ 时，(5.37) 对应两个映射 $\Phi':\boldsymbol{F}\leftrightarrow\boldsymbol{F}_\theta$ 与 $\Phi'':\boldsymbol{F}\leftrightarrow\boldsymbol{F}_\theta$，它们各自交换这两个分支。因此，作用于带平坦丛 $\mathcal{L}$、支撑在 $\boldsymbol{F}$ 上的膜时，$\widetilde{T}_p$ 给出一个支撑在 $\boldsymbol{F}_\theta$ 上、带平坦丛 $(\Phi')^*(\mathcal{L})\oplus(\Phi'')^*(\mathcal{L})$ 的膜；把 $\boldsymbol{F}$ 与 $\boldsymbol{F}_\theta$ 互换时类似。这些都是 $(B,A,A)$ 型膜，可用与我们已经描述过的论证类似的论证来分析。若只考虑支撑在单个分支上的膜，这导致与 (5.35) 对偶的公式；若构造支撑在两分支上的通常磁本征膜，则导致与 (5.36) 对偶的公式。

\subsubsection{improper 't Hooft 算子}

这一切都没有给出几何内窥的好例子，因为 $SL_2$ 局部系统不可能具有不单纯约化到 $SL_2$ 中心的有限自同构群。然而，若我们让 $\widetilde{T}_p$ 作用在并非 $\mathcal{M}_H(SO_3)$ 自身、而是其覆盖的 A-模型上，就能看到内窥在起作用。我们回忆，$\mathcal{M}_H(SO_3,0)$ 与 $\mathcal{M}_H(SO_3,\theta)$ 有覆盖（覆盖群为有限阿贝尔群 $Q=H^1(C,\mathbb{Z}_2)$），它们分别是 $SL_2$ 模空间的 proper 分支与 improper 分支
% ==== 原文 PDF 第 75 页（印刷页 75）====
% [跨页接续: 原文第 74/75 页，句子在此跨页]
$\mathcal{M}_H(SL_2)$ 与 $\mathcal{M}_H(SL_2^*;p)$。
% [背景补充: proper 与 improper 是 [KW] 对 $SL_2$ 平坦丛模空间两个分支的称呼，与「正常态射」无关；按本节末段，它们分别参数化 $\det E=\mathcal{O}$ 与 $\det E=\mathcal{O}(p)$ 的 Higgs 丛，即对应 $w_2$ 为零与非零的 $SO_3$ 丛；译者按（背景补充，无外部依据）。]

improper 分支依赖于基点 $p\in C$ 的选取（如 §5.2.3 所解释）；方便的是取它为我们即将施加 't Hooft 算子的那个点 $p$。

在提升到这些覆盖空间之后，'t Hooft 算子 $\widetilde{T}_p$ 的作用与我们已经描述的相仿。从 $\mathcal{M}_H(SL_2)$ 到 $\mathcal{M}_H(SL_2^*;p)$ 的映射把谱曲线上的线丛与 $\mathcal{O}(p')\oplus\mathcal{O}(p'')$ 作张量积；从 $\mathcal{M}_H(SL_2^*;p)$ 到 $\mathcal{M}_H(SL_2)$ 的映射则把该线丛与 $\pi^*(\mathcal{O}(-p))\otimes(\mathcal{O}(p')\oplus\mathcal{O}(p''))=\mathcal{O}(-p')\oplus\mathcal{O}(-p'')$ 作张量积。这些公式与如下事实相容：两个分支分别参数化 $\det E=\mathcal{O}$ 与 $\det E=\mathcal{O}(p)$ 的 Higgs 丛 $(E,\varphi)$。为确保这一点，其中一个公式中才需要前因子 $\pi^*(\mathcal{O}(-p))$。

% [Tutor Note: 推导补足。$\pi^*(\mathcal{O}(-p))=\mathcal{O}(-p')\otimes\mathcal{O}(-p'')$，因为 $\pi^{-1}(p)=\{p',p''\}$；乘以 $\mathcal{O}(p')\oplus\mathcal{O}(p'')$ 后即得 $\mathcal{O}(-p'')\oplus\mathcal{O}(-p')$，与原文最后一式一致。这里「与 $\mathcal{O}(p')\oplus\mathcal{O}(p'')$ 作张量积」应理解为把线丛的推前分解到目标分支的两个分量上（结果不再是线丛，而是秩 2 的对象），这正是 (5.38) 中 $\widetilde{T}_p\cdot\mathcal{A}_1$ 有两个分量的来源。前因子 $\pi^*(\mathcal{O}(-p))$ 则是沿 improper 到 proper 的方向把行列式的次数从 $\mathcal{O}(p)$ 一侧调回 $\mathcal{O}$ 一侧所必需的；来源：原书 p.75 上下文（A）。]

只要谱曲线光滑，$\widetilde{T}_p$ 在 $SL_2$ 的 A-膜上的作用（交换两个分支）就与我们先前对 $SO_3$ 所描述的类似。但如今我们可以考虑 Hitchin 纤维有两个不可约分支的情形。记 $\boldsymbol{F}_1$ 与 $\boldsymbol{F}_2$ 为 $\mathcal{M}_H(SL_2)$ 的特殊 Hitchin 纤维的两个分支，$\boldsymbol{F}_1^*$ 与 $\boldsymbol{F}_2^*$ 为 $\mathcal{M}_H(SL_2^*;p)$ 的特殊 Hitchin 纤维的两个分支。我们同样记 $\mathcal{A}_1,\mathcal{A}_2$ 与 $\mathcal{A}_1^*,\mathcal{A}_2^*$ 为支撑在这些纤维上的通常类型的 A-膜。

在 't Hooft 算子的作用中，把谱覆盖上的线丛与 $\mathcal{O}(p')$ 作张量积、以及与 $\mathcal{O}(p'')$ 作张量积这两个运算相差一个与 $\mathcal{O}(p'-p'')$ 的张量积。这正是交换两个分支的基本运算。我们可以给点作标号，使得 $\mathcal{O}(p')$ 把 $\boldsymbol{F}_1$ 映到 $\boldsymbol{F}_1^*$、把 $\boldsymbol{F}_2$ 映到 $\boldsymbol{F}_2^*$，而 $\mathcal{O}(p'')$ 把 $\boldsymbol{F}_1$ 映到 $\boldsymbol{F}_2^*$、把 $\boldsymbol{F}_2$ 映到 $\boldsymbol{F}_1^*$。（与 $\mathcal{O}(p')^{-1}$ 或 $\mathcal{O}(p'')^{-1}$ 作张量积当然是逆运算。）没有自然的方式说明哪个是哪个，因为 $p'$ 与 $p''$ 在关于 $p$ 的单值下互换，而这一单值同样互换 $\boldsymbol{F}_1^*$ 与 $\boldsymbol{F}_2^*$，正如我们在 §5.2.3 末所注意到的。

't Hooft 算子 $\widetilde{T}_p$ 在膜 $\mathcal{A}_{1,2}$ 与 $\mathcal{A}_{1,2}^*$ 上的作用示意性地为
\begin{equation}\tag{5.38}
\begin{aligned}
\widetilde{T}_p\cdot\mathcal{A}_1 &= \mathcal{A}_1^*+\mathcal{A}_2^*\\
\widetilde{T}_p\cdot\mathcal{A}_2 &= \mathcal{A}_1^*+\mathcal{A}_2^*,
\end{aligned}
\end{equation}
而把 $\mathcal{A}_i$ 与 $\mathcal{A}_i^*$ 互换时有类似的公式。这些公式以及对偶于三维表示的 't Hooft 算子 $T_p$ 作用的类似公式，与如下关系相容：
\begin{equation}\tag{5.39}
\widetilde{T}_p^2=1+T_p.
\end{equation}
这一关系对偶于式 (5.34)。

% [Tutor Note: 物理图像。(5.39) 与 (5.34) 是同一条张量积关系 $2\otimes2=1\oplus3$ 在两侧的体现：$SO_3$ 虽没有二维表示，但其平方分解中的平凡表示使 $\widetilde{T}_p^2$ 仍是 $SO_3$ 侧合法的算子，从而 (5.39) 才可能成立。$\widetilde{T}_p$ 在 $\mathcal{M}_H(SO_3)$ 上完全自然（对偶于 ${}^{L}G=SL_2$ 的二维表示），但在 $\mathcal{M}_H(SL_2)$ 上只能把膜送到依赖 $p$ 的 $\mathcal{M}_H(SL_2^*;p)$，因为 $SO_3$ 没有二维表示——这就是「improper」的含义；来源：原书 p.75 上下文（A）。]

最后让我们讨论算子 $\widetilde{T}_p$ 有多自然。作为作用在 $\mathcal{M}_H(SO_3)$ 上的膜的算子，它完全自然，因为它对偶于对偶群 ${}^{L}G=SL_2$ 的二维表示。然而，作为作用在 $\mathcal{M}_H(SL_2)$ 上的膜的运算，我们不能指望 $\widetilde{T}_p$ 完全自然，因为它本应对偶于二维表示，而二维表示作为 ${}^{L}G=SO_3$ 的表示并不存在。这种不自然性表现为：若我们希望 $\widetilde{T}_p$ 作用在 $\mathcal{M}_H(SL_2)$ 上的膜上，它把这些膜映到 $\mathcal{M}_H(SL_2^*;p)$ 上的膜，而后者的定义依赖于 $p$（尽管依赖得相对弱）。相比之下，对任何约化群 $G$，与 ${}^{L}G$ 的表示相联系的 't Hooft
% ==== 原文 PDF 第 76 页（印刷页 76）====
% [跨页接续: 原文第 75/76 页，句子在此跨页]
算子，在 $\mathcal{M}_H(G)$ 上的膜上总有完全自然的作用。

作为 $SL_2$ 模空间上的膜的算子，我们称 $\widetilde{T}_p$ 为 improper 't Hooft 算子。
